% v19 complete obstruction. Canonical labels A=0,B=1,C=2,D=3.
% The physical three-high vertex is denoted J in {1,2,3}.
\section{Diamond の四価頂点に一葉を付けた support}
\label{dl19:section}
本節の support は $G_{5,9}=\{01,02,03,04,12,13\}$ である．
葉 $4$ の一高標準形に合わせて残る標識を交換するため，以下では四価頂点を $0$，
三価頂点を $J\in\{1,2,3\}$ と書く．標準形の $(A,B,C,D)$ は $(0,1,2,3)$ であり，
$J$ 以外の二つの標識を $E,F$ とする．従って $I=\{0,1,2,3\}$ の高辺は
$01,02,03,JE,JF$，唯一の低辺は $EF$ である．$J=1,2,3$ の全配置を残す．

一高標準形（命題~\ref{tl15:one-high-full}）は
\[
 Q=q_{04}=th(h+1)+1,\quad u=th,\quad v=t(h+1),\quad h,t\ge1,\quad cd=0,
\]
\[
 V=P^4=\begin{pmatrix}
0&1-Q&-h-1&h\\-1&0&0&0\\-u&u-1-c&0&c\\v&-v-1-d&d&0
\end{pmatrix}.
\]
ここでは $c,d$ は任意の整数であり，$Q=3$ の標準標識の重複も除かない．
$R_{ij}=p^i_{4j}$，$X_{ij}=p^i_{j4}$，$N_{ij}=n_{i4;j}$，
$\ell_0=Q$，$\ell_1=\ell_2=\ell_3=1$ とする．四つの hub 辺の node 条件と $D$ 相殺は
\begin{align}
 R_{ij}&=\ell_i(1-N_{ij})-V_{ij},\quad N_{ij}\ge0,\quad\sum_{j\ne i}N_{ij}=5,
 \label{dl19:reconstruct}\\
 R_{ij}X_{ij}&=\ell_iq_{ij}-1,\qquad
 \sum_{j\ne i}X_{ij}+\sum_{j\ne i}V_{ji}=-2\ell_i.
 \label{dl19:pairD}
\end{align}
core $J$ の低 row は一正二負である（定理~\ref{mh18:matching}）．
cores $E,F$ では，定理~\ref{ls11:two-high-signs} と補題~\ref{ml12:dichotomy} により，
低 row から二つの高 rays へは一正一負，低低 entry は非正である．
低低 entry は一方が零，他方が $-1$ または $-3$ 以下であり，$-2$ は起こらない．
非零の低低 entry を持つ row の高方向 entries は $-S,T$ で，$S\ge2$，$T\ge S$ である．
これらは二高 core の高隣接頂点間の辺が高である場合にも成立する．
従って本節では負の枝 $S=T$ を除かない．

\begin{lemma}[標準 $B$ の接続]
\label{dl19:B-types}
$r=R_{10}$ とすると，$J=1$ では $r=1,2$，$J\ne1$ では $r=1,2,-2$ である．
$r>0$ の場合の $(N_{01},N_{10},q_{01},\kappa)$ は補題~\ref{lh18:B-types} の
四型 I--IV に限られる．ここで $\kappa=1-n_{01;4}\in[-4,1]$ である．
型 I では $Q=3$，型 IV では $Q=3,5$ である．
$r=-2$ の場合にも $Q=3$ である．
\end{lemma}
\begin{proof}
$r=2-N_{10}\in[-3,2]$ は非零である．$J=1$ なら他の二 entries は
$1-N_{12},1-N_{13}\le1$ であり，一正二負と row sum $-1$ より，$r<0$ は $r=-1$ のみである．
$J\ne1$ なら，もう一つの高方向 entry を $s=1-N_{1J}$，低低 entry を $z$ とすると，
$r<0$ の場合には $s=1$，$z=-r-2\le0$ なので $r=-1,-2$ である．
$r=-1$ を除く factorization は補題~\ref{lh18:B-types} の証明と同一であり，
三高 core の仮定を用いない．$r=1,2$ の四型の計算も同様である．
型 I で $Q=3$ となる証明（補題~\ref{lh18:typeI-small}）が用いるのは，
$0B,0C,0D$ が高辺であること，式~\eqref{dl19:pairD} の $i=0$ と辺 $0D$ の node 条件だけである．
従って本 support にも適用できる．

$r=-2$ を調べる．$q=q_{01}=2x+1$，$x\ge1$，$a=2-N_{01}\in[-3,2]$ とすると
\[
 R_{01}=aQ-1,\quad X_{10}=-x,\quad X_{01}=(2\kappa+1)x+\kappa,
\]
\begin{equation}\label{dl19:negative-two-factor}
 (aQ-1)\{(2\kappa+1)x+\kappa\}=Q(2x+1)-1.
\end{equation}
$\kappa=1$ なら $a=1,2$ である．$a=1$ は $(Q-3)x=0$ を与え，
$a=2$ では左辺が右辺を $(4Q-3)x+Q>0$ だけ超える．
$\kappa=0$ でも $a=1,2$ であり，それぞれ
$(Q+1)x+Q-1=0$，$x=1-Q$ を要求して不可能である．
$\kappa=-k<0$ では $1\le k\le4$，$a=-A$，$0\le A\le3$ である．
$A\ge1,k\ge2$ または $A\ge2,k=1$ では左辺が右辺より大きい．
$A=k=1$ では $(Q-1)x=2$ なので $Q=3,x=1$ である．
$A=0$ では $e=Q-k$ と置くと $(2e+1)x=1-e$ であり，
$e\ge1$ と $e\le-1$ は符号が矛盾する．$e=0$ は $Q=k=3,x=1$ を与える．
従って常に $Q=3$ である．
\end{proof}

\begin{proposition}[全ての標準配置における葉分母の上界]
\label{dl19:leaf-bound}
$G_{5,9}$ を support とする五次の算術的接着データが存在するなら $Q\le13$ である．
\end{proposition}
\begin{proof}
$Q>13$ と仮定する．$h\le u$ と $Q=u(h+1)+1$ より，$u\le3$ なら $Q\le3(3+1)+1=13$ となる．従って $u\ge4$，$v\ge5$ である．
前補題により $B$ の接続は型 II または III，$r=R_{10}>0$ に限られる．
常に
\[
 \begin{gathered}
 R_{30}=1-v-N_{30}<0,\\
 R_{21}=2-u+c-N_{21},\qquad R_{23}=1-c-N_{23},\\
 R_{31}=v+2+d-N_{31},\qquad R_{32}=1-d-N_{32}.
 \end{gathered}
\]

$J=1$ の場合，型 II/III と一正二負により $R_{12},R_{13}\in\{-1,-2\}$ である．
core $3$ は二高なので $R_{31}>0$，$R_{32}\le0$ である．
高辺 $13$ の共通低 puncture $4$ における reverse partner の分類
（補題~\ref{ls11:low-puncture}）により $R_{31}=-R_{13}\le2$ である．
しかし row sum は
\[
 v-2+N_{30}=R_{31}+R_{32}\le2
\]
を与え，$v\ge5$ に反する．

以下 $J\ne1$ とする．core $1$ は二高であり，$r=1$ の場合には
$(R_{1J},R_{1E})=(-2,0)$ だけが可能である．
$r=2$ の場合には $(-3,0)$ または $(-2,-1)$ だけが可能である．
実際，低低 entry が $-2$ になる可能性を除き，低低 entry が $-1$ の場合には
二つの高方向 entries の絶対値は同じで二以上でなければならない．
従って常に $R_{1J}\in\{-2,-3\}$ である．

$J=2$ の場合，二高 core $3$ の低低 entry は $R_{31}\le0$ なので
$d\le3-v<0$，従って $c=0$ となる．すると $R_{21}=2-u-N_{21}\le-2$ である．
高辺 $12$ の共通低 puncture において $R_{12},R_{21}\le-2$ となるが，
同符号の負の reverse partners には必ず $-1$ が含まれるので矛盾する．

$J=3$ の場合，高辺 $13$ の reverse partner 分類から
\[
 R_{31}\in\{-1,2,3\}.
\]
$d=0$ なら $R_{31}\ge v-3\ge2$ であり，三高 core $3$ の一正二負条件は
$R_{32}\le-1$ を与える．従って
$v-2+N_{30}=R_{31}+R_{32}\le2$ となり矛盾する．よって $c=0$ である．
$R_{31}>0$ なら同じ row-sum 矛盾になるため，$R_{31}=-1$，
$R_{32}=v-1+N_{30}\ge4$ である．core $2$ は二高で，
$R_{20}=u+1-N_{20}\ge u-4\ge0$ は高方向 entry なので非零，従って正である．
よって $R_{23}=1-N_{23}<0$ であり，辺 $23$ の partner 分類から
$R_{23}=-R_{32}\ge-4$ である．
等号しか残らず，$v=5$，$N_{30}=0$，$N_{23}=5$ となる．
$v=t(h+1)=5$ から $(h,t,u,Q)=(4,1,4,21)$ であり，$N_{20}=N_{21}=0$ なので
$R_{21}=2-u=-2$ となる．これは core $2$ の低低 entry であり，二分補題に反する．
全配置で矛盾したので $Q\le13$ である．
\end{proof}

\subsection{無限の分母族を保持する完全な型検査}

上界 $Q\le13$ によって $(h,t)$ は
\[
 (1,1),\ldots,(1,6),\qquad (2,1),(2,2),\qquad (3,1)
\]
の九組になる．ここから他の分母や $c,d$ が有限であるとは結論しない．
実際，以下の検査では無限の型が現れ，その全体を処理する必要がある．

\begin{lemma}[hub の必要条件の厳密な型表示]
\label{dl19:finite-types}
上の九組，各 $J\in\{1,2,3\}$，各 $cd=0$ の枝，および四つの node 行を固定する．
$cd=0$ の枝を $(c,d)=(z,0)$ または $(0,z)$ と書き，$z$ は全整数を動くものとする．
以下の手順は，hub 行の必要な符号条件，hub の node 条件，および
$I$ 内の高辺の puncture $4$ における pair/clutching 条件を満たす可能性を全て含む．
\begin{enumerate}
\item 各 $0i$ について $R=R_{0i}$，$r=R_{i0}$ とし，$Rr\ne0$ を要求する．
各 $\kappa\in\{-4,-3,-2,-1,0,1\}$ に対し
\[
 \{Qr+R-\kappa Rr\}q_{0i}=R+r
\]
を解く．左係数が非零なら唯一の有理値について，奇数 $q_{0i}\ge3$，
$r\mid q_{0i}-1$，$R\mid Qq_{0i}-1$ を検査する．両辺の係数がともに零なら
\[
 q_{0i}=1+\operatorname{lcm}(2,|r|)a_{0i},\qquad a_{0i}\in\ZZ_{>0}
\]
の全無限族を保持する．
\item $JE,JF$ では，逆方向 entries $r,s$ の組を
\[
 s=-r,\qquad\text{または}\qquad
 \{r,s\}=\{\epsilon,\epsilon(q-1)\},\quad\epsilon\in\{1,-1\}
\]
に限る．通常の $s=-r$ では $q=1+\operatorname{lcm}(2,|r|)a$，$a\ge1$ の全族を残す．
例外の場合には $q=1+\max(|r|,|s|)$ である．
\item core $J$ の低 row に一正二負を課す．core $E,F$ では，高方向 entries が
異符号非零，低低 entry が非正であることを課す．低低 entry が非零の場合，
その row の高方向 entries は $-S,T$ と書け，$S\ge2$，$T\ge S$，$T\ne S+1$ を課す．
これらは必要条件であり，他の row の全局所実現を主張するものではない．
\end{enumerate}
固定された node 行に対し，以上の $z$ の集合は，有限個の点と，偶奇を指定した
完全な半直線の和として厳密に計算できる．全型のうち残るのは，点九百八型と半直線十八型である．
\end{lemma}

\begin{proof}
最初の式は $X_{0i}=(Qq_{0i}-1)/R$，$X_{i0}=(q_{0i}-1)/r$ と
$X_{0i}+X_{i0}=\kappa q_{0i}$ の分母を払ったものである．
零係数の場合は $R=-r$，$r\mid Q-1$ であり，指定した等差数列が全ての整除条件を尽くす．
二番目は共通低 puncture の補題である．固定した $r\ne0$ の許される相手 $s$ は
\[
 \{-r\}\ \cup\
 \begin{cases}\{\operatorname{sgn}(r)\},&|r|\text{ が正の偶数},\\
 \varnothing,&\text{それ以外}\end{cases}
 \ \cup\
 \begin{cases}\{2a\operatorname{sgn}(r):a\ge1\},&|r|=1,\\
 \varnothing,&|r|\ne1\end{cases}
\]
である．可変 row の二 entries は $z+b,-z+b'$ なので，これらの集合の交差は
点と偶奇付き半直線の演算で計算できる．
二高 row の条件も同様である．特に $R_{i0}=r>0$ なら他の高 entry $s$ は
$-r-1\le s\le-1$ の有限区間にあり，$r=-S<0$ なら $S\ge2$ とともに
\[
 s\in\{S-1,S\}\cup\{S+2,S+3,\ldots\}
\]
に限られる．これは低低 entry が零か負かの全ての場合を含む．
従って任意の打切りを導入せずに，有限な node 行と集合演算だけで全型を尽くせる．

付属の \texttt{verify\_g5\_9\_obstruction\_v19.py} はこの手順を整数演算で直接実装する．
四行には各二十一通りを全て使い，九組の $(h,t)$ と三配置を全て検査する．
$z=0$ は $cd=0$ の二つの枝に重複して現れるが，不存在判定に影響しない．
表示した型数はこの標識付き計数であり，同型類の個数ではない．
\end{proof}

\begin{lemma}[残る十八の半直線型は不可能]
\label{dl19:halflines}
補題~\ref{dl19:finite-types} の十八の半直線型は，全て正の二高局所模型の必要条件に反する．
\end{lemma}
\begin{proof}
十八型は全て $J=2$，$(c,d)=(0,z)$ であり，
\[
 (R_{10},R_{12},R_{13})=(1,-2,0),\qquad R_{21}=-1
\]
を満たす．共通低 puncture の例外接着から $q_{12}=3$ となる．
core $1$ の ray $4$ は二高標準形の零低低 row $C$ であり，その高方向 entries が
$1,-2$ なので，canonical な高 rays は $A=0,B=2$ である．
従って canonical な第二高分母 $b$ が三でなければならない．

しかし正の二高標準形で $b=3$ は不可能である．実際，$r+1\mid b-1$ より $r=1$，
$x=a-1$，$y=1$ であり，$u=-s<0$，$s\mid a-1$ から
\[
 a\le A=2a-1-s\le2a-2,\qquad B=(3a-1)/A<3.
\]
$B=1$ は $A=3a-1$ と矛盾し，$B=2$ は $v=B-b+y=0$ を与える．
いずれも標準形の条件に反する．従って十八の半直線の全点が排除される．
\end{proof}

\begin{theorem}[四価頂点に葉を持つ diamond support の不存在]
\label{dl19:no-G5-9}
五次の算術的接着データの高分母 support が $G_{5,9}=\{01,02,03,04,12,13\}$ となることはない．
葉以外の分母にも任意の上限を仮定しない．
\end{theorem}
\begin{proof}
命題~\ref{dl19:leaf-bound}，補題~\ref{dl19:finite-types}，補題~\ref{dl19:halflines} により，
九百八の点型だけを残せばよい．ここで点型とは $z$ の値が定まったことを意味し，
他の分母の自由な等差数列はそのまま残っている．

唯一の低辺 $EF$ については，二高模型の低低 pair の一方が零，他方が負である．
従って $R_{EF}<0$ なら $X_{EF}=0$，$R_{EF}=0$ なら $X_{EF}<0$ であり，
後者は $-2$ にはならない．同じことが $F$ 側にも成り立つ．
低辺 $EF$ の puncture $4$ における node 条件は
\[
 -4\le X_{EF}+X_{FE}\le1
\]
である．両 entries は非正なので，$R$ が零である側の $X$ は
$\{-1,-3,-4\}$ のいずれかに限られる．以上から $X_{EF},X_{FE}$ を有限に選べる．

他の五辺について，自由な分母をそれぞれ $q_{ij}=1+L_{ij}a_{ij}$，$a_{ij}\ge1$ と書く．
全 $X_{ij}$ はこの自由変数の整数係数の affine 式として復元される．
式~\eqref{dl19:pairD} の四本の $D$ 相殺は，有理数係数の線形系になる．
以下は三配置と二つの $cd=0$ 枝ごとの全計数である．可変 row 二は $c$，三は $d$ を表す．
\[
\begin{array}{cc|rr|rr}
 J&\text{可変 row}&z\text{ 点型}&z\text{ 半直線型}&D\text{ 不整合}&D\text{ 正値性矛盾}\\\hline
 1&2&69&0&179&4\\
 1&3&84&0&227&5\\
 2&2&108&0&321&3\\
 2&3&121&18&338&5\\
 3&2&142&0&398&4\\
 3&3&384&0&1100&12\\\hline
 \multicolumn{2}{c|}{\text{合計}}&908&18&2563&33
\end{array}
\]
低辺の $X$ の複数の選択により，点型九百八個から線形系二千五百九十六個が生じる．
二千五百六十三個では有理係数の行基本変形が $0=\gamma\ne0$ を与える．
残る三十三個は全て $Q=3$，$h=t=1$ であり，row $R_{0*}$ は順序を除き $(-7,2,2)$，
負 entry に対応する分母は五，残りの二分母は $1+2a,1+2b$ である．
$\sum_jV_{j0}=-1-u+v=0$ なので，$i=0$ の $D$ 相殺は
\[
 -2+(1+3a)+(1+3b)=-6,
 \qquad\text{すなわち}\qquad 3a+3b=-6
\]
を要求する．$a,b\ge1$ に反する．

付属検算は全ての線形系を \texttt{fractions.Fraction} によって処理し，
各矛盾について元の四式の有理線形結合を証明書に記録する．
さらに，記録する線形結合を元の係数行列へ再乗算して検査する．
自由変数には数値上限を課していないので，これは有限探索範囲内の不存在ではなく，
残る全分母族の排除である．
\end{proof}

本定理が排除するのは，固定された標準 $K_5$ branch graph と marked punctures を持つ
全平面五次 SNC 構成の当該 support である．一般型曲面上の任意の Lagrangian torus fibration
に対する非存在定理ではない．一高の任意整数 $c,d$，分母三の標準形の重複，二高の負の枝，
および他の分母の全自由族をいずれも保持している．
